\pdfoutput=1
\documentclass[11pt]{article}

\usepackage[margin=1in]{geometry}
\usepackage{amsmath,amssymb}
\usepackage{booktabs}
\usepackage{microtype}
\usepackage{graphicx}
\usepackage{url}
\usepackage[hidelinks]{hyperref}
\usepackage{natbib}
\usepackage{xcolor}

\title{Forward-Only Adapter Fine-Tuning at 0.9B Parameters:\\
A Systems Extension of Dust-Inspired Activation Perturbation}
\author{Independent technical report\\
Authors and affiliations to be finalized before submission}
\date{Draft: October 2026}

\newcommand{\dust}{\textsc{Dust}}
\newcommand{\kmodel}{K2-Horizon-0.9B}
\newcommand{\oproj}{\texttt{o\_proj}}

\begin{document}
\maketitle

\begin{abstract}
Dust recently demonstrated that activation-space zeroth-order optimization can
train transformer language models without a backward pass, with gradient
alignment to backpropagation improving as population grows
\citep{dahal2026dust}. We study a complementary setting: adapting a
\emph{pretrained} 0.9B-parameter decoder-only transformer using a rank-4
low-rank adapter on the final attention output projection, while keeping all
pretrained weights frozen. Our experiments are an independent
\emph{Dust-inspired systems extension}, not a reproduction of the complete
Dust pretraining algorithm.

We first show that a naive one-sided Gaussian estimator can be too noisy for
this setting. Antithetic orthogonal perturbations improve alignment
substantially; a complete 1536-direction basis yields local LoRA-B gradient
cosine 0.981 against exact autograd, and after two real adapter updates the
resulting A/B update cosines reach 0.998/0.991. We then exploit the restricted
fine-tuning site: because the adapter is attached to the final decoder
layer's attention \oproj{}, the frozen prefix can be cached and only the
downstream residual/MLP/norm/language-model-head tail replayed for each
perturbation. This reduces a fixed four-step K=1024 three-seed protocol from
347.4 seconds to 15.0 seconds on average, a 23.2$\times$ speedup over our
serial structured implementation and approximately 2.1$\times$ the elapsed
time of a matched adapter-only backpropagation control.

The quality result remains negative/neutral: average heldout cross-entropy
worsens slightly for both forward-only and backprop controls in the tiny
protocol. We also implement an independent per-token Gaussian estimator closer
to Dust's tokenwise perturbation semantics. It is viable but does not beat the
shared orthogonal control at K=256--1024 in this position-local final-layer
setting. Our results therefore support a narrower claim: forward-only
activation perturbation can recover backprop-like adapter updates at this
scale, and systems structure can remove much of the apparent execution cost,
but update fidelity is not sufficient evidence of generalization.
\end{abstract}

\section{Introduction}

Backpropagation is the dominant credit-assignment mechanism for modern neural
networks, but it imposes differentiability and backward-graph requirements on
the computation being optimized. Zeroth-order methods replace analytic
gradients with loss evaluations under controlled perturbations, potentially
making a broader class of systems trainable at the cost of additional forward
compute \citep{salimans2017es,nesterov2017random,malladi2023zo}.

\dust{} advances this direction by perturbing hidden activations rather than
weights. Its central observation is that independent token-level activation
perturbations can act as many virtual population members inside a transformer
forward pass. The reported pretraining experiments show that larger
populations improve gradient alignment and that activation-space
zeroth-order training can approach or exceed backpropagation in selected
settings, albeit with substantially more compute \citep{dahal2026dust}.
The accompanying implementation exposes populations from 256 to 16,384 and
uses multi-GPU runs for the larger settings \citep{qlabsdustcode}.

We ask a different question: \emph{can the same broad idea be useful when
adapting a pretrained language model rather than pretraining every weight?}
Low-rank adaptation (LoRA) restricts fine-tuning to small trainable matrices
inserted into an otherwise frozen model \citep{hu2021lora}. This creates a
particularly favorable setting for forward-only learning: a perturbation
estimator need only recover the local update for a small adapter, while the
pretrained network remains fixed.

Our study uses \kmodel{}, a 28-layer pretrained decoder-only transformer with
hidden width 1536, and attaches a rank-4 LoRA to the final decoder layer's
attention \oproj{}. We compare exact adapter-only autograd with several
activation-perturbation estimators, measure update cosine directly, and then
optimize the systems path by caching the frozen prefix.

The main contributions of this draft are:
\begin{itemize}
    \item An independent evaluation of Dust-inspired activation-space
    zeroth-order adapter tuning on a pretrained 0.9B-parameter transformer,
    including AMD ROCm execution.
    \item A gradient-calibration study showing that estimator fidelity rises
    sharply with structured antithetic population size, reaching near-autograd
    update direction at a complete hidden-space basis.
    \item A cached-tail execution strategy that preserves the constrained
    final-\oproj{} estimator while making the fixed K=1024 protocol roughly
    23$\times$ faster than our initial serial implementation.
    \item A deliberately negative quality result: high update cosine does not
    yet translate into repeatable heldout improvement in our small protocol.
    \item A completed independent per-token Gaussian perturbation
    comparison showing that Dust-like tokenwise noise is viable here but is
    less sample-efficient than the shared orthogonal control at K=256--1024,
    clarifying where our fine-tuning setting differs from the full method.
\end{itemize}

\paragraph{Scope.}
This paper is not a reproduction of the full Dust pretraining result. We do
not perturb every linear layer, do not implement Dust's attention
future-token credit assignment, and do not claim that our orthogonal
shared-direction estimator is identical to Dust. The current result is best
understood as an independent systems and fine-tuning extension motivated by
the same activation-space zeroth-order principle.

\section{Related Work}

\paragraph{Dust.}
\citet{dahal2026dust} perturb linear-layer outputs with Gaussian noise,
independently at every token, and use token-local loss changes to estimate
error signals. Attention internals receive specialized credit that includes
future-token effects. Their experiments pretrain GPT-style transformers on
FineWeb, report three seeds per setting, and study how population, model size,
and token budget affect test loss and gradient alignment. The public code is a
minimal implementation that preserves the estimator and tuned defaults while
omitting some execution optimizations \citep{qlabsdustcode}.

\paragraph{Evolution strategies and zeroth-order optimization.}
Weight-space evolutionary strategies estimate gradients from loss differences
under random parameter perturbations \citep{salimans2017es}. Structured and
orthogonal directions can improve sampling efficiency and variance
\citep{choromanski2018structured}. Random gradient-free optimization also has
a classical theory connecting finite-difference estimates to smoothed
objectives \citep{nesterov2017random}. More recently, forward-only
fine-tuning has been studied directly for language models
\citep{malladi2023zo}.

\paragraph{Low-rank adaptation.}
LoRA represents a parameter update to a weight matrix through low-rank factors,
reducing trainable state while preserving the frozen pretrained weights
\citep{hu2021lora}. Our setting uses LoRA not as a claim about the optimal
fine-tuning method, but as a controlled substrate for directly measuring
forward-only update fidelity.

\section{Method}

\subsection{Model and adapter}

We freeze all pretrained parameters of \kmodel{}. The model has 28 decoder
layers and hidden width $d=1536$. A rank-$r=4$ LoRA is attached only to the
final decoder layer's attention output projection $W_o$.

For input activation $x$, the adapted projection is
\begin{equation}
    h = x W_o^\top + (x A^\top) B^\top,
\end{equation}
where $A \in \mathbb{R}^{r\times d}$ and
$B \in \mathbb{R}^{d\times r}$. The pretrained $W_o$ remains unchanged.

\subsection{Forward-only local gradient estimator}

For hidden output $h$, perturbation direction $u_k$, noise scale $\sigma$,
and supervised loss $L$, our antithetic estimator is
\begin{equation}
\hat g_h =
\frac{1}{2\sigma K}
\sum_{k=1}^{K}
\left[
L(h+\sigma u_k)-L(h-\sigma u_k)
\right] u_k.
\label{eq:antithetic}
\end{equation}

The estimated local error is mapped into LoRA gradients using the captured
\oproj{} input. With $z=xA^\top$,
\begin{align}
    \widehat{\nabla_B L} &= \hat g_h^\top z, \\
    \widehat{\nabla_A L} &= (\hat g_h B)^\top x.
\end{align}
No backward pass is invoked in this path.

Our initial experiments used Gaussian directions. We then use randomized
orthogonal directions scaled by $\sqrt{d}$ so their second moment is
comparable to isotropic Gaussian perturbations. When $K=d=1536$, the
directions form a complete randomized orthogonal basis.

\subsection{Difference from Dust}

This distinction is central. Dust samples activation perturbations
independently at each token position \citep{dahal2026dust}. Our current
structured K2 estimator shares each direction across token positions within a
draw. That choice made the gradient-calibration experiment simple and made a
complete hidden-space basis possible, but it does not realize Dust's
token-as-virtual-population mechanism.

We therefore add an independent per-token antithetic Gaussian estimator
and compare its compute/fidelity frontier directly with the shared-direction
control. The tokenwise estimator is viable but less sample-efficient at
K=256--1024 in this final-layer, position-local setting. We retain the term
Dust-inspired because earlier-layer perturbations and attention future-token
credit remain outside the present experiment.

\subsection{Cached final-tail replay}

Because the adapter is attached only to the final attention \oproj{}, all
computation upstream of that projection is invariant to the temporary LoRA
factors and perturbations. We evaluate the frozen prefix once per example and
cache the final decoder-layer input, the final \oproj{} input, and the clean
final \oproj{} output.

Every perturbation then replays only
\begin{equation}
\text{modified }\oproj{}
\rightarrow \text{residual}
\rightarrow \text{final MLP}
\rightarrow \text{final norm}
\rightarrow \text{LM head}
\rightarrow L.
\end{equation}

Masked prompt positions are excluded from scored-tail replay. This
optimization is exact only because the trainable/perturbed site is the final
\oproj{}. It cannot be applied unchanged to earlier layers or q/k/v
projections, where perturbations affect future attention states.

\section{Experimental Setup}

\paragraph{Hardware and software.}
Primary GPU experiments ran on an AMD Radeon 8050S through PyTorch
2.12.0+rocm7.14.0. The model was loaded in BF16 for the ROCm experiments.
The base weight shard was pinned by SHA-256 for every run.

\paragraph{Data.}
Examples were selected deterministically by SHA-256 from an existing cleaned
local corpus. Heldout selection rejects exact normalized prompt overlap with
the training source. Raw prompts, responses, and token IDs are not serialized
into public result manifests. The small quality protocol uses 16 training
examples (864 supervised assistant tokens) and 12 prompt-disjoint heldout
examples (624 supervised tokens).

\paragraph{Controls.}
The principal control is a fresh rank-4 LoRA with identical initialization,
trained by exact \texttt{torch.autograd.grad} on the adapter factors only.
All base-model parameters remain frozen. Reported comparisons use the same
example order and fixed selection digests.

\paragraph{Integrity.}
Research runs do not merge or save trained adapters. Base weights are checked
for equality, base gradients must remain absent, and the forward-only path
reports zero backward calls. Immutable result manifests include source/model
digests and are retained separately from private training text.

\section{Results}

\subsection{Gradient calibration}

The first real-model pilot used one-sided Gaussian perturbations with
$K=64$ and $\sigma=0.05$. On the calibration example, its LoRA-B gradient
cosine against exact local autograd was only approximately 0.106. This
motivated antithetic and structured perturbations.

Table~\ref{tab:gradient} reports the antithetic orthogonal sweep at
$\sigma=0.25$.

\begin{table}[t]
\centering
\caption{Local gradient alignment as orthogonal population grows.}
\label{tab:gradient}
\begin{tabular}{rccc}
\toprule
Population $K$ & Output-grad cosine & LoRA-B cosine & LoRA-B rel. $\ell_2$ \\
\midrule
256  & 0.40162 & 0.40411 & 2.2518 \\
512  & 0.56927 & 0.54898 & 1.4131 \\
1024 & 0.79906 & 0.79675 & 0.7435 \\
1536 & \textbf{0.98563} & \textbf{0.98125} & \textbf{0.1947} \\
\bottomrule
\end{tabular}
\end{table}

The full 1536-direction basis required 3074 model forwards and approximately
121.9 seconds in the original full-model calibration. The result shows that
the local forward-only estimator can closely recover the analytic gradient
direction when given sufficient structured population.

\subsection{Multi-step update fidelity}

We next compare actual adapter trajectories, not a single local gradient.
After two updates, both $A$ and $B$ can move because $B$ is no longer zero
after the first step.

\begin{table}[t]
\centering
\caption{Two-step adapter update cosine versus matched backprop.}
\label{tab:updates}
\begin{tabular}{rccr}
\toprule
$K$ & LoRA-A cosine & LoRA-B cosine & Forward-only time (s) \\
\midrule
512  & 0.32584 & 0.60441 & 88.16 \\
1024 & 0.93034 & 0.81621 & 177.29 \\
1536 & \textbf{0.99796} & \textbf{0.99144} & 262.32 \\
\bottomrule
\end{tabular}
\end{table}

The complete hidden-space basis nearly reproduces the matched backprop update
after two real adapter updates. A separate four-step $K=1024$ run retains
LoRA-A cosine 0.9936 and LoRA-B cosine 0.8552, showing that estimator
coherence can persist after both factors have changed.

\subsection{Systems optimization}

Full-model microbatching reduces Python/model-call overhead but changes
BF16/ROCm arithmetic slightly because repeated identical sequences are
evaluated in a larger batch. Caching a single-example frozen prefix both
reduces compute and avoids this source of estimator drift.

For the formal $K=1024$, two-update comparison, the serial full-model path
takes 154.24 seconds. Cached scored-tail replay takes 6.36 seconds for the
update loop plus 0.80 seconds to cache the prefixes, while preserving A/B
update cosine 0.9997/0.9996 versus the serial reference.

\begin{table}[t]
\centering
\caption{K=1024 two-step systems comparison.}
\label{tab:systems}
\begin{tabular}{lrrrr}
\toprule
Method & Update time (s) & Prefix (s) & A cosine & B cosine \\
\midrule
Serial full model & 154.24 & -- & reference & reference \\
Cached scored-tail & \textbf{6.36} & 0.80 & \textbf{0.9997} & \textbf{0.9996} \\
\bottomrule
\end{tabular}
\end{table}

This is a 24.3$\times$ update-loop speedup and 21.5$\times$ including prefix
construction.

\subsection{Three-seed quality protocol}

We then repeat the fixed four-step $K=1024$ protocol over seeds
7, 42, and 1337.

\begin{table}[t]
\centering
\caption{Four-step fixed quality protocol. Negative CE delta is improvement.}
\label{tab:quality}
\begin{tabular}{rrrrrr}
\toprule
Seed & Tail train $\Delta$ & Tail heldout $\Delta$ & Tail time & Serial time & BP time \\
\midrule
7    & -0.000011 & +0.000813 & 14.81 & 346.89 & 7.46 \\
42   & -0.000858 & +0.001127 & 15.31 & 347.57 & 6.87 \\
1337 & +0.000407 & -0.000177 & 14.87 & 347.80 & 6.86 \\
\midrule
Mean & -0.000154 & +0.000588 & \textbf{14.99} & 347.42 & 7.06 \\
\bottomrule
\end{tabular}
\end{table}

The cached forward-only path is 23.2$\times$ faster than the earlier serial
structured implementation and approximately 2.12$\times$ matched backprop
for this tiny protocol. The mean train and heldout deltas closely track the
serial structured reference, which is the expected behavior for a
semantics-preserving systems optimization.

The quality conclusion is intentionally conservative. One forward-only seed
improves the tiny heldout set, but average heldout CE still worsens. Matched
backprop also worsens average heldout CE in this protocol. These data do not
support a claim of superior generalization.

\section{Discussion}

\subsection{What the results support}

The strongest evidence is about \emph{credit assignment and execution}, not
downstream model quality. A large structured perturbation population can
recover a local adapter update that is nearly collinear with exact autograd.
The effect survives multiple updates. When the trainable site is localized
near the output, a systems implementation can exploit frozen-prefix structure
to remove most of the redundant compute.

These findings are consistent with Dust's broader argument that
activation-space zeroth-order estimators can remain meaningful at transformer
scale \citep{dahal2026dust}. They add evidence in a different regime:
pretrained-model adapter tuning, a 0.9B-parameter base model, and AMD ROCm.

\subsection{What the results do not support}

We do not show that forward-only fine-tuning beats backpropagation. We do not
show that our current shared-direction estimator is the same as Dust. We do
not test attention future-credit, earlier layers, or full-model parameter
updates. Our heldout set is intentionally small and is unsuitable for a broad
generalization claim.

The distinction matters because a method can produce an update highly aligned
with backprop and still fail to improve task quality under a short schedule.
In our current experiments, that is exactly what occurs.

\subsection{Why the negative result is useful}

A contribution to the Dust line of work is more credible if it reports where
the forward-only estimator works and where the evidence stops. Our experiments
suggest that the basic estimator is not obviously blocked by model scale or
by AMD execution. Instead, the remaining research problem is a combination of
population efficiency, estimator semantics, and optimization schedule.

\section{Tokenwise Perturbation Parity}

To close the largest semantic gap with Dust, we implemented independent
per-token Gaussian activation perturbations in the cached final-\oproj{} tail.
For scored token $t$ and draw $k$, the estimator uses an independently sampled
direction $u_{k,t}$:
\begin{equation}
\hat g_{h_t} =
\frac{1}{2\sigma K}
\sum_{k=1}^{K}
\left[
L_t(h_t+\sigma u_{k,t})-
L_t(h_t-\sigma u_{k,t})
\right]u_{k,t}.
\end{equation}

We compare this tokenwise Gaussian estimator with the shared orthogonal
antithetic control over seeds 7, 42, and 1337, populations
$K\in\{64,256,1024\}$, and $\sigma\in\{0.05,0.25\}$. The calibration
state uses a reproducible small nonzero LoRA-B factor so both A and B gradients
are measurable.

At the preferred $\sigma=0.25$, mean cosine to exact local autograd is:

\begin{table}[t]
\centering
\caption{Three-seed tokenwise-parity calibration at $\sigma=0.25$.}
\label{tab:tokenwise}
\begin{tabular}{llrrrr}
\toprule
Estimator & $K$ & Output cosine & LoRA-A cosine & LoRA-B cosine & Time (s) \\
\midrule
Shared orthogonal & 64   & 0.1964 & 0.1984 & 0.1899 & 0.117 \\
Tokenwise Gaussian & 64 & 0.1976 & 0.2059 & 0.2033 & 0.116 \\
Shared orthogonal & 256  & \textbf{0.4002} & \textbf{0.3401} & \textbf{0.4075} & 0.467 \\
Tokenwise Gaussian & 256 & 0.3716 & 0.2344 & 0.3780 & 0.466 \\
Shared orthogonal & 1024 & \textbf{0.8091} & \textbf{0.7278} & \textbf{0.8116} & 1.875 \\
Tokenwise Gaussian & 1024 & 0.6266 & 0.5024 & 0.6274 & 1.868 \\
\bottomrule
\end{tabular}
\end{table}

The tokenwise estimator is viable and improves with population, but it does
not outperform the shared orthogonal estimator at the operating populations.
At $K=1024$, the shared estimator has materially higher output and LoRA
gradient cosine at essentially identical measured tail runtime.

This negative result is informative rather than contradictory to Dust. In our
restricted final-\oproj{} experiment, the downstream tail is position-local:
each scored token's loss depends only on its own perturbed final-layer output.
Reusing the same well-structured orthogonal basis across token positions does
not create cross-token credit contamination in this setting, while
orthogonality reduces estimator variance. Dust addresses a broader case in
which many sites are perturbed and attention creates future-token credit.

A two-step real-model training smoke verifies that the selectable tokenwise
path is operational end-to-end with zero backward calls and an unchanged base
model. We therefore treat tokenwise parity as a completed calibration gate,
not as a reason to replace the current structured estimator.

\section{Next Experiment: Quality Scaling}

Estimator selection is now frozen for the next quality experiment:
shared orthogonal antithetic perturbations with $K=1024$, $\sigma=0.25$, and
D=4 cached scored-tail replay.

The next protocol will test 8, 16, and 32 update schedules across multiple
predeclared seeds, with a larger independently checked heldout set,
near-duplicate screening, fixed tool-call exactness/JSON-validity metrics, and
fixed agentic completion tasks. The matched adapter-only backprop control will
remain in every cell.

The estimator will not be retuned using heldout results. Checkpoint promotion
requires repeatable heldout or task-level improvement; train-loss reduction or
high gradient cosine alone remains insufficient.

\section{Limitations}

\begin{itemize}
    \item The adapter modifies only the final attention \oproj{}. This is a
    deliberately narrow setting.
    \item Current structured perturbations are shared across token positions,
    unlike Dust's independent per-token activation noise.
    \item The cached-tail optimization depends on the trainable site being at
    the final layer and cannot be generalized mechanically to earlier layers.
    \item The quality protocol is small, and exact prompt-disjointness does
    not guarantee semantic non-overlap.
    \item Only a small number of seeds and one pretrained base model have been
    studied.
    \item Wall-time comparisons are systems-specific and should not be read
    as universal hardware ratios.
    \item We compare to adapter-only backpropagation, not full-model
    backpropagation.
\end{itemize}

\section{Reproducibility and Release Plan}

The intended public release contains estimator code, synthetic fixtures,
manifest schemas, aggregate result tables, and pinned model/code digests.
Private prompts, responses, token IDs, credentials, and machine-specific
secrets are excluded.

We intend to share the independent results with the Dust authors as a
systems/fine-tuning extension. Any authorship or joint-paper arrangement
should be based on direct collaboration and contribution rather than implied
by citation or upstream inspiration.

\section{Conclusion}

Dust motivates a compelling question: how far can forward-only
activation-space credit assignment scale before the cost of zeroth-order
estimation becomes prohibitive? In a constrained but realistic adapter-tuning
setting, we find that the answer depends strongly on implementation structure.
A sufficiently large structured population recovers backprop-like adapter
updates on a 0.9B pretrained transformer, and cached-tail replay removes most
of the redundant full-model compute. The resulting tiny protocol is roughly
2.1$\times$ matched adapter-only backprop rather than tens of times slower.

The remaining open question is more important than execution fidelity:
whether the fixed best estimator, evaluated over longer predeclared schedules
and broader task suites, produces repeatable heldout or task-level gains. The
tokenwise parity gate is now complete and does not justify replacing the
shared orthogonal K=1024 control in this setting.

\bibliographystyle{plainnat}
\bibliography{references}

\appendix
\section{Current result ledger}

\begin{table}[h]
\centering
\caption{Selected immutable results used in this draft.}
\begin{tabular}{ll}
\toprule
Result & Value \\
\midrule
Naive K=64 LoRA-B gradient cosine & $\approx 0.106$ \\
K=1536 local LoRA-B gradient cosine & 0.98125 \\
K=1536 two-step A/B update cosine & 0.99796 / 0.99144 \\
K=1024 cached-tail two-step A/B cosine vs serial & 0.999665 / 0.999577 \\
K=1024 cached-tail two-step update-loop speedup & 24.26$\times$ \\
Three-seed cached-tail mean runtime & 14.99 s \\
Three-seed serial structured mean runtime & 347.42 s \\
Three-seed matched backprop mean runtime & 7.06 s \\
Three-seed cached-tail mean heldout CE delta & +0.000588 \\
Tokenwise K=1024 A/B cosine at $\sigma=.25$ & 0.5024 / 0.6274 \\
Shared orthogonal K=1024 A/B cosine at $\sigma=.25$ & 0.7278 / 0.8116 \\
\bottomrule
\end{tabular}
\end{table}

\section{Submission checklist}

Before any arXiv submission:
\begin{enumerate}
    \item finalize authors and affiliations;
    \item rerun all tables from immutable public-safe manifests;
    \item regenerate tokenwise-parity tables from the immutable public-safe manifest;
    \item expand heldout/task evaluation;
    \item verify every hardware/runtime claim against manifests;
    \item replace provisional web citations if an official Dust arXiv entry
    becomes available;
    \item add a code/data availability statement with the public release URL;
    \item perform an independent technical read for overclaiming.
\end{enumerate}

\end{document}
